\documentclass[11pt]{article}
\usepackage[T1]{fontenc}
\usepackage{lmodern}
\usepackage[margin=1in]{geometry}
\usepackage{amsmath,amssymb,amsthm,mathtools}
\usepackage[colorlinks=true,linkcolor=blue,urlcolor=blue,citecolor=blue]{hyperref}
\newtheorem{theorem}{Theorem}
\newtheorem{corollary}[theorem]{Corollary}
\newtheorem{remark}[theorem]{Remark}
\DeclareMathOperator{\QR}{QR}
\DeclareMathOperator{\QNR}{QNR}

\title{A ten-point deficit progression in the mod-107\\
signed-divisor support}
\author{Erd\H{o}s--Straus Workbench\\
community lead credited to \texttt{u/CommonCareful3149}}
\date{25 August 2026}

\begin{document}
\maketitle

\begin{abstract}
For the record integer $p=8{,}803{,}369$ and residual shell $R=107$, the
even-factor half-log support has a ten-point complement in
$\mathbb Z/53\mathbb Z$.  We verify that this complement is the arithmetic
progression $D=\{23+42j:0\leq j\leq9\}$ and that $|D-D|=19$.  Since the gap
between the two odd-sheet translates is not in $D-D$, the translates saturate
the full quadratic-nonresidue coset modulo $107$.  This is one bounded support
calculation, not a general shell theorem, an endpoint proof, or a priority
claim.
\end{abstract}

\section{Provenance and scope}

A public post by Reddit user
\href{https://www.reddit.com/r/LLMmathematics/comments/1vy127z/}
{u/CommonCareful3149} reports an investigation of the mod-$107$
Erd\H{o}s--Straus shell and announces a short technical note.  This workbench
credits the progression and difference-set observation to that community lead and
reconstructs it independently below.  The unpublished note and private
correspondence are not reproduced.

The exact progression language was not found in the routed local project
corpus; the factor logarithms, parity sheets, target fibres, and missing
exponent classes were already present there.  This local search result does
not establish global novelty.  The immutable historical baseline is
\href{https://zenodo.org/records/21845035}{Erd\H{o}s--Straus Project Archive,
Zenodo record 21845035}.

\section{Coordinates}

Set
\[
p=8{,}803{,}369,
\qquad R=107,
\qquad a=\frac{p+R}{4}=2{,}200{,}869=3^2 11^2 43\,47.
\]
The element $2$ has order $106$ modulo $107$.  Hence
\[
\lambda:(\mathbb F_{107}^{\times},\cdot)\longrightarrow
 C_{106}:=\mathbb Z/106\mathbb Z,
\qquad \lambda(u)=\log_2(u),
\]
is an isomorphism, with inverse $e\mapsto2^e$.  Direct computation gives
\begin{equation}
 \lambda(3)=70,\qquad \lambda(11)=22,\qquad
 \lambda(43)=59,\qquad \lambda(47)=66.
 \label{eq:logs}
\end{equation}
Thus $43$ is the unique factor with odd logarithm.

The quadratic residues correspond under $\lambda$ to the even subgroup
$2C_{106}$.  Define the exact half-log isomorphism
\[
h:2C_{106}\longrightarrow C_{53}:=\mathbb Z/53\mathbb Z,
\qquad h(2x)=x,
\]
with inverse $x\mapsto2x$.  Removing the factor $43$, the centered
even-factor support maps to
\begin{equation}
A=\{35r+11s+33t:-2\leq r,s\leq2,\ -1\leq t\leq1\}
\subseteq C_{53}.
\label{eq:A}
\end{equation}
The map $h$ has trivial kernel and singleton fibres; this change of coordinate
loses no information.

\section{The deficit and its difference set}

\begin{theorem}
The set $A$ has $43$ elements.  Its complement is
\begin{align}
D:=C_{53}\setminus A
 &=\{1,10,12,21,23,30,32,41,43,52\}\notag\\
 &=\{23+42j\pmod{53}:0\leq j\leq9\}.
\label{eq:D}
\end{align}
Moreover,
\begin{equation}
D-D=\{42k\pmod{53}:-9\leq k\leq9\},
\qquad |D-D|=19=2|D|-1.
\label{eq:DD}
\end{equation}
\end{theorem}

\begin{proof}
Exact reduction of the $5\cdot5\cdot3=75$ triples in \eqref{eq:A} gives
$|A|=43$ and the first list in \eqref{eq:D}.  Starting at $23$ and repeatedly
adding $42$ gives
\[
23,12,1,43,32,21,10,52,41,30,
\]
which proves the progression identity.

Every difference is $42(j-i)$ with $-9\leq j-i\leq9$, and each integer in
that interval occurs.  Because $42$ is invertible modulo $53$, a collision
would imply $53\mid(k-\ell)$.  But $|k-\ell|\leq18<53$, so $k=\ell$.
Therefore all $19$ displayed differences are distinct.
\end{proof}

\section{Saturation of the nonresidue sheet}

The odd exponent coset is a torsor, not a subgroup.  We therefore use the
affine bijection
\[
\psi:1+2C_{106}\longrightarrow C_{53},
\qquad \psi(2x+1)=x,
\qquad \psi^{-1}(x)=2x+1.
\]
Since
\[
59=2\cdot29+1,
\qquad -59\equiv47=2\cdot23+1\pmod{106},
\]
the two nonzero centered choices for the $43$ exponent give supports
$A+29$ and $A+23$.  Their complements are $D+29$ and $D+23$.

\begin{theorem}
The two translated deficits are disjoint, and hence
\[
(A+29)\cup(A+23)=C_{53}.
\]
Under $\psi^{-1}$ and exponentiation by $2$, the two centered $43$-factor
choices saturate the entire quadratic-nonresidue coset in
$\mathbb F_{107}^{\times}$.
\end{theorem}

\begin{proof}
An intersection would give $d_1+29=d_2+23$, so $6=d_2-d_1\in D-D$.  But
\[
42^{-1}=24\pmod{53},
\qquad24\cdot6=144\equiv38\equiv-15\pmod{53}.
\]
The integer $-15$ is not in $[-9,9]$, so \eqref{eq:DD} implies
$6\notin D-D$, a contradiction.  Taking complements proves the union
identity.  Both subsequent coordinate maps are bijections.
\end{proof}

In the alternative Legendre-sheet coordinate $v\mapsto h(\lambda(-v))$, the
same supports are $A+3$ and $A-3$.  The two affine presentations differ by a
translation and have the same decisive gap $6$.

\begin{corollary}
The full centered exponent support
\[
B=\{70r+22s+59u+66t:
-2\leq r,s\leq2,\ -1\leq u,t\leq1\}\subseteq C_{106}
\]
has $43$ even and $53$ odd classes.  Thus $|B|=96$ and
\[
C_{106}\setminus B=2D
=\{2,20,24,42,46,60,64,82,86,104\}.
\]
The target exponent $53=\lambda(-1)$ has exactly the two coefficient fibres
\[
(-2,0,-1,-1),\qquad(2,0,1,1),
\]
whereas the exterior target exponent $60=\lambda(-p^{-1})$ has none.
\end{corollary}

\begin{proof}
Only $59u$ controls parity.  The case $u=0$ gives the $43$-point even support;
$u=\pm1$ gives the saturated $53$-point odd support.  The complement is the
inverse half-log image $2D$.  Direct enumeration of the finite coefficient
box gives the two stated fibres and the empty exterior fibre.
\end{proof}

\section{Reproducibility and nonclaims}

The standard-library checker
\begin{center}
\texttt{certificates/verify\_r107\_deficit\_progression.py}
\end{center}
reconstructs the primitive-generator table, both support coordinates, $D$,
$D-D$, both translate presentations, residue-level nonresidue saturation,
the full missing set, and both target fibres.  Ordinary and optimized Python
runs must have byte-identical output; the file contains no removable
\texttt{assert} nodes.

The Lean file
\begin{center}
\texttt{formal/lean/R107DeficitProgression.lean}
\end{center}
independently checks the finite-set equalities in $\operatorname{ZMod}(53)$
and $\operatorname{ZMod}(106)$ and the factor logarithms in
$\operatorname{ZMod}(107)$.  Formal acceptance and correspondence with this
informal argument are separate review obligations.

\begin{remark}
Support saturation does not by itself retain representation multiplicities;
the fibre enumeration is a separate calculation.  Nothing here proves a
family of shells, a density theorem, or the Erd\H{o}s--Straus conjecture.
The workbench makes no claim that the progression observation is new.
\end{remark}

\end{document}
